\chapter{Introdução}
\label{chap:intro}
\begin{itemize}
	
	\item INTRODUZIR SOBRE SIMULAÇÃO ( PARA QUE SERVE)
	\item SIMULAÇÃO HIBRIDA ( PARA QUE SERVE)
	\item Gêmeos digitais Atrelado a simulação hibrida
	\item Aplicaçao: Construção e citar que há problemas
	\item CONTEXTO: Minha casa minha vida
	
	
\end{itemize}


\section{Justificativa}
\label{sec:justificativa}
\begin{itemize}
	\item Justificar por que a solução hibrida é uma solução valida
	\item (lacuna)Muito tradicional/ Manual/ Boca a boca
	\item (lacuna)Dizer que os estudos não tem uma parte visual -> omniverse;
	\item dizer o por que é importante/necessário;
	\item Ir alimentando
	
	
\end{itemize}


Em alguns pontos, a construção civil ainda se mantém desatualizada. Segundo \citeonline{cheng}, na gestão tradicional da construção civil as informações relevantes de canteiros de obras frequentemente comunicadas por meio de documentos em papel ou comunicação verbal.

\section{Objetivos}
\label{sec:objetivos}
O objetivo geral deste trabalho é propor uma simulação híbrida para gestão e controle de recursos físicos do sistema construtivo de paredes de concreto, visando otimizar o fluxo de trabalho no canteiro de obras e aumentar a produtividade.
Para alcançar este propósito, definem-se os seguintes objetivos específicos:
\begin{itemize}

	\item Desenvolver um modelo computacional de simulação híbrida utilizando a ferramenta AnyLogic;
	\item Criar a representação visual contextualizada em um Gêmeo Digital do canteiro através do software NVIDIA Omniverse;
	\item Integrar os dados de simulação ao ambiente virtual para análise de cenários.
\end{itemize}

\section{Organização do trabalho}
\label{sec:organizacao}
Este trabalho está estruturado em cinco capítulos, organizados da seguinte forma:

O Capítulo \ref{chap:intro} apresenta a contextualização do tema, a justificativa para a realização da pesquisa, bem como os objetivos propostos.

O Capítulo \ref{chap:revisao} dedica-se à revisão da literatura, abordando os conceitos fundamentais sobre Gêmeos Digitais, Simulação Híbrida e as ferramentas AnyLogic e Omniverse.

O Capítulo \ref{chap:metodologia} descreve a metodologia da pesquisa, detalhando as etapas de construção do modelo conceitual, coleta de dados e a implementação computacional.

O Capítulo \ref{chap:resultados} expõe a validação do modelo, a análise dos cenários simulados 

O Capítulo \ref{chap:discussao} faz a discussão sobre os impactos da integração das ferramentas na gestão do concreto.

Por fim, o Capítulo \ref{chap:conclusao} apresenta as considerações finais, as limitações encontradas ao longo do estudo e as sugestões para trabalhos futuros.






